\chapter{Organisation de la Cellule de Crise}
\label{chap:cellulecrise}

La première victime d'une crise est souvent la structure hiérarchique habituelle. Pour réagir vite, l'entreprise doit se structurer en une organisation temporaire, centralisée et souple : la Cellule de Crise.

\section{Présentation : La Structure de Commandement}
% ------------------------------------------------------------------------------

\subsection{Le Principe de Centralisation}

En temps de crise, on ne peut pas avoir plusieurs centres de décision qui se contredisent. Il faut un commandement unique. L'organisation se divise généralement en trois niveaux.

\begin{figure}[htbp]
\centering
\begin{tikzpicture}[
    node distance=1.5cm,
    box/.style={rectangle, draw, thick, fill=white, minimum width=4cm, minimum height=1cm, align=center}
]
    % Nodes
    \node[box, fill=red!20] (strat) {Cellule Stratégique\\(Direction)};
    \node[box, fill=orange!20, below left=of strat] (com) {Cellule Communication\\(Image)};
    \node[box, fill=blue!20, below right=of strat] (op) {Cellule Opérationnelle\\(Technique)};
    
    % Arrows
    \draw[->, thick] (strat) -- node[left, font=\scriptsize] {Consignes} (com);
    \draw[->, thick] (strat) -- node[right, font=\scriptsize] {Consignes} (op);
    \draw[->, thick, dashed] (op) -- node[below, font=\scriptsize] {Retours terrain} (strat);
    \draw[->, thick, dashed] (com) -- node[below, font=\scriptsize] {Retours médias} (strat);
\end{tikzpicture}
\caption{Organisation type d'une Cellule de Crise}
\end{figure}

\subsection{Les Trois Niveaux}

\begin{enumerate}
    \item \textbf{Cellule Stratégique :} Le cerveau. Elle prend les décisions majeures (Arrêter l'activité ? Payer la rançon ? Demander la mise en liquidation ?). Elle est composée de la Direction Générale.
    \item \textbf{Cellule Opérationnelle :} Les mains. Elle exécute les décisions sur le terrain (Technique, Logistique, Sécurité Physique). Elle remonte l'information factuelle.
    \item \textbf{Cellule Communication :} La voix. Elle gère l'image et la diffusion de l'information (Presse, Clients, Réseaux Sociaux).
\end{enumerate}

% ------------------------------------------------------------------------------
\section{Rôles et Responsabilités (Jargon)}
% ------------------------------------------------------------------------------

\subsection{Le Directeur de Crise}

\begin{boxdef}
\textbf{Directeur de Crise :} La personne qui anime la cellule de crise et tranche les arbitrages. Il a l'autorité ultime sur les ressources de l'entreprise pour la durée de la crise.
\end{boxdef}

C'est généralement le PDG ou le Directeur Général. Il doit être calme, décideur et protecteur de l'intérêt général de l'entreprise.

\subsection{Le Porte-parole (Spokesperson)}

\begin{boxdef}
\textbf{Porte-parole :} Seule personne habilitée à s'exprimer officiellement auprès des médias et des parties prenantes externes.
\end{boxdef}

\begin{boxalert}
\textbf{Règle d'Or :} Une seule voix ! Personne d'autre ne doit s'exprimer. Si un employé parle à la presse sans autorisation, il peut désinformé et aggraver la crise.
\end{boxalert}

\begin{boxlizbiz}
\textbf{Dilemme LiZbiZ :}
Le PDG est très charismatique mais panique facilement. La DSI est très technique mais austère.
\textbf{Choix stratégique :} Le PDG assume le rôle de Porte-parole pour rassurer (Leadership), mais il est briefé en amont par la DSI sur les faits techniques précis.
\end{boxlizbiz}

\subsection{Le Scribe (Secrétaire de Crise)}

Rôle souvent sous-estimé mais vital.

\begin{boxdef}
\textbf{Scribe :} Personne chargée de noter par écrit la chronologie des événements, les décisions prises, les actions lancées et les horaires exacts.
\end{boxdef}

Le cahier du Scribe constitue la preuve de la diligence de l'entreprise ("Nous avons réagi à 09h05"). C'est une pièce maîtresse en cas de procès futur.

\subsection{La Salle de Crise}

Lieu physique sécurisé où se réunit la cellule stratégique.
\begin{itemize}
    \item Doit être équipée (Téléphones, Projecteur, Tableaux, Alimentation électrique de secours).
    \item Doit être isolée (Confidentialité des débats).
    \item Peut être un "PC Mobile" si les locaux sont inaccessibles (ex: incendie).
\end{itemize}

% ------------------------------------------------------------------------------
\section{Cas LiZbiZ : Mise en Place de l'Organigramme}
% ------------------------------------------------------------------------------

\subsection{Scénario : Activation de la Cellule}

\begin{boxlizbiz}
\textbf{Situation :} La fuite de données est confirmée. Le PDG décide d'activer le Plan de Gestion de Crise à 09h00.
\end{boxlizbiz}

\subsection{Travaux Pratiques : Constitution des Équipes}

Les étudiants doivent répartir les rôles suivants au sein de la classe (ou en sous-groupes) :

\begin{table}[htbp]
\centering
\begin{tabular}{p{3cm}p{5cm}p{4cm}}
    \toprule
    \textbf{Rôle} & \textbf{Profil / Compétence} & \textbf{Mission Immédiate} \\
    \midrule
    Directeur de Crise & PDG (Décideur). & Arbitrer la stratégie globale. \\
    Porte-parole & Dir. Com. ou PDG. & Préparer le communiqué initial. \\
    Scribe & Juriste ou Assistant. & Noter l'heure de chaque événement. \\
    Animateur Technique & DSI / RSSI. & Faire le point sur l'étendue de la fuite. \\
    \bottomrule
\end{tabular}
\caption{Tableau de répartition des rôles LiZbiZ}
\end{table}

\subsection{Exercice : La Première Réunion}

\textbf{Consigne :} Simulez les 15 premières minutes de la réunion de crise.

\begin{enumerate}
    \item Le Directeur de Crise ouvre la séance.
    \item La DSI fait un point situation (Ce qu'on sait, ce qu'on ne sait pas).
    \item Le Scribe note : "09h05 - Confirmation fuite données santé. 09h07 - Décision d'appeler l'avocat."
    \item Débat : Faut-il prévenir la CNIL tout de suite ou attendre d'en savoir plus ? (Réponse légale : Oui, délai de 72h).
\end{enumerate}

% ==============================================================================
% Fichier : 05_module_crise.tex (Extrait Chapitre 3)
% Description : Communication de Crise
% ==============================================================================